\documentclass{article}
\usepackage[T1]{fontenc}
\usepackage{lmodern}
\usepackage{graphicx}
\usepackage{amsmath,amssymb}
\usepackage[a4paper,margin=2cm]{geometry}
\usepackage[absolute,overlay]{textpos}
\setlength{\TPHorizModule}{1mm}
\setlength{\TPVertModule}{1mm}
\usepackage[indonesian]{babel}
\usepackage{hyperref}
\setlength{\emergencystretch}{2em}
% unit: Halaman 69

\begin{document}

% === Halaman 69 (python/native) ===

• Sifat reversible ini membuat perancang cipher blok tidak perlu 

membuat algoritma baru untuk mendekripsi cipherteks menjadi 

plainteks. 

 Contoh: 

Li – 1  f(Ri – 1, Ki)  f(Ri – 1, Ki) = Li – 1 

• Sifat reversible tidak bergantung pada fungsi f sehingga fungsi f dapat 

dibuat serumit mungkin.

69

\end{document}
